\cleardoublepage

\begin{frontmatter}
	\chapter{Capítulo de Exemplo da segunda parte\footnote{This is chapter footnote}}\label{chap3}
	\subchapter{Chapter Subtitle}

	\begin{aug}
		\author[addressrefs={ad1,ad2}]%
		{\fnm{Firstname}   \snm{Surname}}%
		%{\fnm{Firstname}   \snm{Surname}\footnote{This is author footnote}}%
		\author[addressrefs={ad2}]%
		{\fnm{Firstname} \snm{Surname}}%
		\address[id=ad1]%
		{Short Address}%
		\address[id=ad2]%
		{Long Address}%
	\end{aug}

	% \minitoc

	%
	\begin{chapterpoints}%[Chapter Points]
		\item The ends of words and sentences are marked by spaces. It doesn't
		matter how many spaces you type; one is as good as 100.  The end of
		a line counts as a space.

		\item The ends of words and sentences are marked by spaces. It doesn't
		matter how many spaces you type; one is as good as 100.  The end of
		a line counts as a space.
	\end{chapterpoints}

	\begin{dispquote}

		The ends of words and sentences are marked by spaces. It doesn't
		matter how many spaces you type; one is as good as 100.  The end of
		a line counts as a space.

		The ends of words and sentences are marked by spaces. It doesn't
		matter how many spaces you type; one is as good as 100.  The end of
		a line counts as a space.

		\source{Name}

	\end{dispquote}

\end{frontmatter}

\section{Section title}\label{sec3.1}

Texto para tirar o erro de citação nula \citep{lopes2024}

\section{Código Fortran - Exemplo}

\begin{lstlisting}
PROGRAM ola
    PRINT *, "Olá, Mundo!"
END PROGRAM ola
\end{lstlisting}

\lstinputlisting{src/test.f90}
